%% Section 01 — Title page and candidature type
\begin{titlepage}
  \begin{center}
    \vspace*{2cm}

    {\Large Queensland University of Technology}\\[0.5em]
    {\large Faculty of Science}\\[0.5em]
    {\large School of Mathematical Sciences}

    \vspace{2.5cm}

    {\LARGE\bfseries
      Computational Modelling of Multilateral Treaties\\
      and Consensus Dynamics:\\[0.4em]
      AI, Game Theory, and the Antarctic Treaty System
    }

    \vspace{2cm}

    \textbf{Stage 2 — Confirmation of Candidature}

    \vspace{1.5cm}

    \begin{tabular}{rl}
      \textbf{Candidate:}        & Oriolson Rodriguez \\[4pt]
      \textbf{Student Number:}   & n12608521 \\[4pt]
      \textbf{Degree:}           & Doctor of Philosophy (PhD) \\[4pt]
      \textbf{Mode of Study:}    & Full-time \\[4pt]
      \textbf{Submission Date:}  & \today \\
    \end{tabular}

    \vfill

    {\small
      Submitted in partial fulfilment of the requirements\\
      of the degree of Doctor of Philosophy at\\
      Queensland University of Technology.
    }
  \end{center}
\end{titlepage}

%% ── Abstract ────────────────────────────────────────────────────────────────
\chapter*{Abstract}
\addcontentsline{toc}{chapter}{Abstract}

International environmental treaties such as the Antarctic Treaty System (ATS) depend on
consensus decision-making among heterogeneous state actors with divergent interests.
Despite decades of scholarship in international law and political science, the precise
computational mechanisms that drive treaty stability, coalition formation, and consensus
collapse remain poorly understood.

This project develops a formal computational framework — integrating agent-based modelling,
game-theoretic analysis, and machine-learning techniques — to model how multilateral
consensus is reached, maintained, and broken within treaty regimes.
Using the ATS as a primary empirical case, the project will: (\textit{i}) characterise
the affinity and aversion dynamics that govern state negotiating behaviour;
(\textit{ii}) simulate consensus emergence under varying institutional rules; and
(\textit{iii}) identify leverage points for treaty reform and stability.

The research contributes new computational tools and theoretical insights applicable
beyond the ATS to climate, fisheries, and arms-control regimes.

\vspace{1em}
\noindent\textbf{Keywords:} Antarctic Treaty System, consensus dynamics, multilateral treaties,
agent-based modelling, game theory, computational social science, international cooperation.
